\begin{table}[htb]
\centering
\setlength{\tabcolsep}{8pt}
\begin{tabular}{l|c|c}
\hline
Method  & \multicolumn{1}{c|}{AP@IoU=0.5} & \multicolumn{1}{|c}{AP drop} \\ \hline
AttFuse~\cite{xu2022opv2v} & 22.5  & 42.2                       \\
AttFuse \textbf{w/ D.A.}  & 23.4 \color{gray}{(+0.9)} & 41.3                           \\\arrayrulecolor{lightgray}\hline
F-Cooper~\cite{chen2019f} & 23.6  & 37.1                         \\
F-Cooper \textbf{w/ D.A.}   & 37.3 \color{gray}{(+13.7)} & 23.4             \\\arrayrulecolor{lightgray}\hline
V2VNet~\cite{wang2020v2vnet}   & 23.2  &  41.3                        \\
V2VNet \textbf{w/ D.A.}   & 26.3 \color{gray}{(+3.1)}  & 38.2          \\\arrayrulecolor{lightgray}\hline
V2X-ViT~\cite{xu2022v2xvit} & 27.4  &  37.5          \\
V2X-ViT \textbf{w/ D.A.}    & 39.5 \color{gray}{(+12.1)}  &  25.4      \\\arrayrulecolor{lightgray}\hline
CoBEVT~\cite{xu2022cobevt} & 32.6 &  33.9          \\
CoBEVT \textbf{w/ D.A.}   & 40.2 \color{gray}{(+7.6)} &  26.3                          \\ \arrayrulecolor{black}\hline
\end{tabular}
\vspace{-2mm}
\caption{\textbf{Domain Adaptation benchmark}. The number in the bracket indicates the precision gain when using domain adaptation. AP drop refers to the precision gap compared to directly training on the V2V4Real dataset.}
\label{tab:da}
\vspace{-3mm}
\end{table}
